% ============================================================
% ОБЩАЯ-СЕМАНТИКА: Вирусо-логия семантических реакций
% + Дополнение: Школа как фабрика инъекций
%
% SocioCyberneering Inc. — 2026
% XeLaTeX + polyglossia
% Шрифты: Liberation Serif (основной), Montserrat (заголовки/обложка)
% ============================================================
\documentclass[13pt,a4paper]{article}

\usepackage{fontspec}
\usepackage{polyglossia}
\setmainlanguage{russian}
\setotherlanguage{english}

\setmainfont{Liberation Serif}
\setsansfont[Scale=0.98]{Liberation Sans}
\setmonofont{DejaVu Sans Mono}[Scale=0.85]
\newfontfamily\montserrat[Path=./fonts/,Extension=.ttf,UprightFont=*-Regular,BoldFont=*-Bold,ItalicFont=*-Regular,BoldItalicFont=*-Bold]{MontserratStatic}

\usepackage{geometry}
\geometry{margin=2.5cm, top=3cm, bottom=3cm}

\usepackage{xcolor}
\usepackage{hyperref}
\usepackage{enumitem}
\usepackage{booktabs}
\usepackage{array}
\usepackage{tabularx}
\usepackage{parskip}
\usepackage{microtype}
\usepackage{graphicx}
\usepackage{titlesec}
\usepackage{setspace}
\usepackage{nameref}
\usepackage{mdframed}
% ---- схемы в стиле «Доска Мерсера» ----
\usepackage{needspace}
\usepackage{pdfpages}
\usepackage{tikz}
\usetikzlibrary{arrows.meta,positioning,calc,backgrounds,fit,shapes.geometric,decorations.pathreplacing}
\definecolor{board}{HTML}{F6EFDD}
\newcommand{\ttl}{\montserrat\bfseries\footnotesize}
\newcounter{scheme}
\newenvironment{scheme}[2]{%
  \par\Needspace{#2}\vspace{0.9em}%
  \noindent{\montserrat\bfseries\small\color{sce-gold}\refstepcounter{scheme}СХЕМА \thescheme\ \textbullet\ #1}\par
  \vspace{-0.2em}\noindent\textcolor{sce-gold}{\rule{\linewidth}{0.8pt}}\par\vspace{0.4em}\nopagebreak
  \begin{center}\small
}{%
  \end{center}\vspace{0.4em}%
}
\tikzset{
  card/.style={draw=sce-black, line width=1pt, rounded corners=2pt, align=center,
    inner sep=5pt, font=\small, text=sce-black},
  cb/.style={card, fill=bg-blue,  draw=sce-blue},
  cg/.style={card, fill=bg-green, draw=sce-green},
  cr/.style={card, fill=bg-red,   draw=sce-red},
  cy/.style={card, fill=bg-gold,  draw=sce-gold},
  cw/.style={card, fill=white},
  thread/.style={sce-red, line width=1.2pt},
  arr/.style={-{Stealth[length=3mm,width=2.4mm]}, line width=1.4pt, sce-black},
  arrr/.style={-{Stealth[length=3mm,width=2.4mm]}, line width=1.6pt, sce-red},
  arrg/.style={-{Stealth[length=3mm,width=2.4mm]}, line width=1.6pt, sce-green},
  arrd/.style={-{Stealth[length=3mm,width=2.4mm]}, line width=1.2pt, dash pattern=on 4pt off 3pt, sce-grey},
  pin/.style={circle, fill=sce-red, draw=sce-black, line width=0.4pt, inner sep=1.7pt},
}
\newcommand{\boardbg}{%
  \begin{scope}[on background layer]
    \fill[rounded corners=4pt, fill=board, draw=sce-grey, line width=0.6pt]
      ([shift={(-5mm,-5mm)}]current bounding box.south west) rectangle
      ([shift={(5mm,5mm)}]current bounding box.north east);
  \end{scope}}


% ---- глоссарий: термины с якорями и кликабельными ссылками ----
% \term{КОД}{перевод} → [КОД] (перевод) — кликабельно, прыгает в глоссарий
% \termdef{КОД}{перевод} → определяет якорь в глоссарии
\newcommand{\term}[2]{%
  \hyperref[glossary:#1]{%
    \texttt{\textcolor{sce-blue}{[#1]}}%
  }%
  \,{\footnotesize\textcolor{sce-grey}{[#2]}}%
}
% краткий вариант без перевода (когда перевод уже ясен из контекста)
\newcommand{\termr}[1]{%
  \hyperref[glossary:#1]{\textbf{\texttt{\textcolor{sce-blue}{[#1]}}}}%
}
% якорь определения в глоссарии
\newcommand{\termdef}[1]{\phantomsection\label{glossary:#1}}

% красные кавычки: \kav{слово} → 'слово' красным
\newcommand{\kav}[1]{\textcolor{sce-red}{'#1'}}

% ---- цвета SCE ----
\definecolor{sce-blue}{HTML}{154360}      % чернильный синий — механизм/нейтраль
\definecolor{sce-gold}{HTML}{9A7D0A}      % бронзово-золотой — диалог/цитата
\definecolor{sce-red}{HTML}{7B241C}       % тёмно-красный — опасность/предупреждение
\definecolor{sce-grey}{HTML}{2D3748}      % графитовый — вспомогательный
\definecolor{sce-green}{HTML}{145A32}     % тёмно-зелёный — факт/антивирус/решение
\definecolor{sce-black}{HTML}{0D0D0D}     % почти чёрный — основной текст заголовков

\hypersetup{
  colorlinks=true,
  linkcolor=sce-black,
  urlcolor=sce-blue,
}

% ---- стиль заголовков ----
\titleformat{\section}
  {\montserrat\large\bfseries\color{sce-black}\addfontfeatures{LetterSpace=5}}
  {\thesection.}{0.8em}{}[\vspace{-0.3em}\textcolor{sce-gold}{\rule{\linewidth}{0.6pt}}]

\titleformat{\subsection}
  {\montserrat\normalsize\bfseries\color{sce-blue}}
  {\thesubsection.}{0.6em}{}

\titleformat{\subsubsection}
  {\sffamily\small\bfseries\color{sce-grey}}
  {\thesubsubsection.}{0.5em}{}

% ---- фоновые цвета для боксов (светлые) ----
\definecolor{bg-blue}{HTML}{EBF5FB}
\definecolor{bg-gold}{HTML}{FEF9E7}
\definecolor{bg-red}{HTML}{FDEDEC}
\definecolor{bg-green}{HTML}{EAFAF1}

% ---- box-окружения — левая полоса + фон ----
% mechbox — СИНИЙ — механизм, определение, нейтральный анализ
\newenvironment{mechbox}[1]{%
  \vspace{0.9em}
  \begin{mdframed}[
    linecolor=sce-blue, linewidth=4pt,
    topline=false, rightline=false, bottomline=false,
    backgroundcolor=bg-blue,
    innerleftmargin=12pt, innerrightmargin=10pt,
    innertopmargin=8pt, innerbottommargin=8pt,
    skipabove=0pt, skipbelow=0pt
  ]
  \noindent{\montserrat\normalsize\bfseries\color{sce-blue}#1}\par\vspace{4pt}%
  \ignorespaces
}{%
  \end{mdframed}
  \vspace{0.6em}
}

% dialogbox — ЗОЛОТОЙ — цитата, пример из жизни, диалог
\newenvironment{dialogbox}[1]{%
  \vspace{0.9em}
  \begin{mdframed}[
    linecolor=sce-gold, linewidth=4pt,
    topline=false, rightline=false, bottomline=false,
    backgroundcolor=bg-gold,
    innerleftmargin=12pt, innerrightmargin=10pt,
    innertopmargin=8pt, innerbottommargin=8pt,
    skipabove=0pt, skipbelow=0pt
  ]
  \noindent{\montserrat\normalsize\bfseries\color{sce-gold}#1}\par\vspace{4pt}%
  \ignorespaces
}{%
  \end{mdframed}
  \vspace{0.6em}
}

% warnbox — КРАСНЫЙ — опасность, предупреждение, атака
\newenvironment{warnbox}[1]{%
  \vspace{0.9em}
  \begin{mdframed}[
    linecolor=sce-red, linewidth=4pt,
    topline=false, rightline=false, bottomline=false,
    backgroundcolor=bg-red,
    innerleftmargin=12pt, innerrightmargin=10pt,
    innertopmargin=8pt, innerbottommargin=8pt,
    skipabove=0pt, skipbelow=0pt
  ]
  \noindent{\montserrat\normalsize\bfseries\color{sce-red}#1}\par\vspace{4pt}%
  \ignorespaces
}{%
  \end{mdframed}
  \vspace{0.6em}
}

% factbox — ЗЕЛЁНЫЙ — антивирус, инструмент, решение, факт
\newenvironment{factbox}[1]{%
  \vspace{0.9em}
  \begin{mdframed}[
    linecolor=sce-green, linewidth=4pt,
    topline=false, rightline=false, bottomline=false,
    backgroundcolor=bg-green,
    innerleftmargin=12pt, innerrightmargin=10pt,
    innertopmargin=8pt, innerbottommargin=8pt,
    skipabove=0pt, skipbelow=0pt
  ]
  \noindent{\montserrat\normalsize\bfseries\color{sce-green}#1}\par\vspace{4pt}%
  \ignorespaces
}{%
  \end{mdframed}
  \vspace{0.6em}
}

% ============================================================
\begin{document}

% ---- ОБЛОЖКА ----
\thispagestyle{empty}
\begin{titlepage}
\centering
\vspace*{3cm}

% Надстрочная метка
{\montserrat\bfseries\small\addfontfeatures{LetterSpace=12}\color{sce-black}
О Б Щ А Я - С Е М А Н Т И К А}

\vspace{0.4em}
{\montserrat\bfseries\footnotesize\addfontfeatures{LetterSpace=8}\color{sce-black}
Н А У К А  \ Ч Е Л О В Е Ч Е С К И Х  \ О Ц Е Н О К}

\vspace{1.8em}
\textcolor{sce-gold}{\rule{0.4\textwidth}{0.6pt}}
\vspace{1.2em}

% Главный заголовок
{\montserrat\fontsize{32}{38}\selectfont\bfseries\color{sce-black}
Вирусо-логия\\[0.2em]
семантических\\[0.2em]
реакций}

\vspace{1.4em}
\textcolor{sce-gold}{\rule{0.3\textwidth}{0.4pt}}
\vspace{1.6em}

% Подзаголовок
{\montserrat\bfseries\normalsize\addfontfeatures{LetterSpace=4}\color{sce-black}
Дополнение к общей-семантике}

\vspace{0.5em}
{\small\color{sce-black}\itshape
Включает: Вирусология оценок~--- Школа}

\vfill

% Нижний блок — как на обложке книги
\textcolor{sce-gold}{\rule{0.5\textwidth}{0.3pt}}
\vspace{0.8em}

{\montserrat\bfseries\small\addfontfeatures{LetterSpace=8}\color{sce-black}
S O C I O C Y B E R N E E R I N G}

\vspace{0.2em}
{\montserrat\bfseries\footnotesize\addfontfeatures{LetterSpace=5}\color{sce-black}
I N C .}

\vspace{0.6em}
{\small\color{sce-black} 2026 \quad|\quad v\,0.2 \quad|\quad GPL-3.0}

\vspace{0.4em}
{\footnotesize\color{sce-black}
\href{https://sociocyber.site}{sociocyber.site} \quad|\quad
\href{https://git.sociocyber.site}{git.sociocyber.site} \quad|\quad
\href{https://github.com/sociocyber1}{github.com/sociocyber1}}

\end{titlepage}

% ---- ОГЛАВЛЕНИЕ ----
\newpage
\tableofcontents
\newpage

% ============================================================
%  ЧАСТЬ I — Вирусо-логия семантических реакций (оригинал)
% ============================================================

\newpage
\begin{center}
\vspace*{1cm}
\textcolor{sce-gold}{\rule{0.5\textwidth}{0.4pt}}\\[0.8em]
{\montserrat\Large\bfseries\addfontfeatures{LetterSpace=3}\color{sce-black}
ЧАСТЬ I}\\[0.4em]
{\montserrat\large\addfontfeatures{LetterSpace=5}\color{sce-grey}
В И Р У С О - Л О Г И Я}\\[0.2em]
{\montserrat\normalsize\addfontfeatures{LetterSpace=3}\color{sce-grey}
семантических реакций}\\[0.8em]
\textcolor{sce-gold}{\rule{0.5\textwidth}{0.4pt}}
\vspace{1cm}
\end{center}

% ============================================================
\section{Примечание~--- версия 0.1v}
% ============================================================

Данный документ ещё не готов. Опубликован заранее, так как возможно, содержит что-то
важное, на что раньше, насколько я могу судить, не обращали внимания. Вероятно содержит
ошибки (разные): как в оценках, в моих \kav{картах} \kav{карт}, в том как \kav{что-то} я назвал, то что
\kav{пронаблюдал}-\kav{испытал} в той или иной степени. Некоторые разделы будут изменены, добавлены
пояснения, графика, картинки, по возможности реальные примеры (если вы таковые имеете~---
поделитесь, мы добавим).

\subsection{Не используй меня во вред!}

Я не хочу, чтобы вы намеренно и сознательно использовали материалы в злонамеренных,
разрушительных, деструктивных целях, таких как получение выгоды из людей и их затупление
чем-то \kav{до-научным}.

Данный документ~--- пока что, что-то около анализа войн, проблем в наших оценках, что-то
явно не так, и я пытаюсь в этом разобраться.

\subsection{Инъекции от меня~--- берегись!}

Я допустил себя к рассуждениям во введении~--- возможно, они не такие полезные.
Аккуратнее с~--- моей оценкой. Возможно я выбрал не самое полезное название в
разделе \kav{Болезни ИФД}~--- первый термин, кажется, содержит некоторую ошибку.
Вам потребуется рассмотреть вашу проблему с разных точек-областей-наук, учесть не только
нейро-оценочную-среду, но и микробиологию, внешние/внутренние среды.

\subsection{Перепроверь меня!}

Надеюсь, вы не как Шерлок Холмс~--- делаете выводы поспешно, сразу без каких-либо проверок
своих гипотез и доказательств. Следы, улики~--- можно подделать; индукции-дедукции будут
ошибочны. Выводы~--- неверны. Карта~--- подстроена или неверна.

\begin{warnbox}{No jump to conclusions! — Не прыгайте к выводам}
Я дал общую \kav{карту}. Я возможно ошибся.\\
Ваш \kav{Опыт1} и мой \kav{Опыт2} не повторяются в точности.\\
Там, где нет контролируемого эксперимента~--- скорее всего туфта, а не наука.
\end{warnbox}

% ============================================================
\section{Введение}
% ============================================================

Добро пожаловать в \textbf{Вирусо-логию человеческих оценок} и семантических реакций на них!

Данное исследование проводится в рамках проекта SocioCyberneering, а также в рамках
\textit{SocioCyberneering: Architecture Design Pre-Definitions For Safety-Live-Behavior-In-Semantic-Reaction-Environments}.

Массовые несчастья и кое-что похуже привели меня в ужас и побудили искать ответы.
Чтобы привести человеческие отношения в симбиоз с планетой, без войн, нищеты, политики,
требуется выяснить причины самих войн, источники самой \kav{ненависти} и общего безрассудства.

\begin{mechbox}{Ключевые вопросы}
Какие вирусы в нашем \kav{программном биологическом обеспечении} существуют?\\[0.3em]
Как они передаются?\\[0.3em]
\kav{Символы} преобразуют нас \kav{биохимически} и меняют наши \kav{нейро-оценочные транзакции}?\\[0.3em]
Что за \kav{семантические болезни} тянутся к нам из книг прошлого?\\[0.3em]
Что такого они там читают~--- \kav{военные министры}, \kav{короли/президенты}?
\end{mechbox}

Полагаться на символы~--- это как ехать на зелёный свет светофора, но сбивать пешеходов.
Символ сообщил: езжай. А наблюдение~--- упущено. Результат: сбит человек.

Отчёты~--- символы. Книги~--- символы. Старая рукопись~--- символы.
Прямое наблюдение~--- упущено; в прошлое нельзя просто заглянуть и проверить.

\begin{dialogbox}{Пример~--- как символы ведут к войне}
Король записал: \kav{отомстить за людей и покорить Францию}.\\
Это ещё не значит, что кто-то у него погиб.\\
Он мог сделать фальшивую запись, чтобы побудить принца начать войну.\\
Называть по-разному: \kav{там живут одни безбожники и еретики}.\\[0.4em]
Карта заменила территорию.\\
Война~--- следствие.
\end{dialogbox}

% ============================================================
\section{Болезни ИФД}
% ============================================================

В своей книге Уэлдел Джонсон писал:

\begin{dialogbox}{Уэлдел Джонсон~--- о болезни ИФД}
Идеалы в паттерне ИФД стремятся в высь в трёх отношениях: их размыто определяют;
высоко ценят; и в их определении плохо согласуют с реальностью.\\[0.4em]
Неспособность распознать свои достижения приводит человека к чувству фрустрации, которое,
в свою очередь, ведёт к (а)~ощущению некоторой степени неполноценности и (б)~последующим
реакциям, от агрессии до явного уныния и утраты инициативы.\\[0.4em]
Болезнь ИФД: от идеализма~--- к фрустрации~--- и затем к деморализации.\\
Масштабы её распространённости близятся к пандемическим.
\end{dialogbox}

\begin{scheme}{Цикл ИФД по Уэнделлу Джонсону}{7.5cm}
\begin{tikzpicture}
  \node[cy, text width=4.4cm] (i) at (0,2.6)     {\textbf{Идеализм}\\идеал размыт, высок, не согласован с реальностью};
  \node[cr, text width=4.4cm] (f) at (5.2,-1.2)  {\textbf{Фрустрация}\\достижений не видно, чувство неполноценности};
  \node[cr, text width=4.4cm, line width=1.6pt] (d) at (-5.2,-1.2) {\textbf{Деморализация}\\агрессия, уныние, потеря инициативы};
  \draw[arrr] (i) -- (f);
  \draw[arrr] (f) -- (d);
  \draw[arrr] (d) -- node[above left=-1pt,font=\footnotesize\itshape,text=sce-red]{новый, ещё менее достижимый идеал} (i);
  \node[font=\montserrat\bfseries\normalsize, text=sce-red, align=center] at (0,0) {ИФД};
  \foreach \n in {i,f,d} \node[pin] at (\n.north west) {};
  \boardbg
\end{tikzpicture}
\end{scheme}


Волнения, страхи и дурные предчувствия составляют значительную долю наших личностных
дезадаптаций. Мы испытываем \textbf{оценочный страх}: сами факты не пытают нас морально
настолько, насколько мы это делаем своими оценками этих фактов~--- и чертями, которых
создаём из воздуха.

Болезнь ИФД~--- не столько недуг индивидуумов, сколько отражение влияния на них
семантических сил. Мы легко \kav{подхватываем семантические болячки}.

\begin{mechbox}{Как передаются семантические болячки?}
Подсказка~--- в неспособности дезадаптированных людей вербализовать свои трудности
таким образом, чтобы на них получилось ответить эффективно.\\[0.3em]
Это указывает на \textbf{язык}~--- на то, что в структуре нашего общего языка присутствуют
разрушительные факторы.\\[0.3em]
Эту проблему разобрал Альфред Коржибски в книге \textit{«Наука и Здравомыслие»},\\
и~--- в более сжатой форме~--- Сэмюэл Хаякава в книге \textit{«Язык в Действии»}.
\end{mechbox}

% ============================================================
\section{Перехват контроля в обход \kav{коры}. Роль Амигдалы в безумии}
% ============================================================

\begin{mechbox}{\termr{E.A.Alert}~--- Emygdala Alert~--- Тревога амигдалы}
Атака, при которой злоумышленник или случайное событие переводит человека в режим
\kav{бей-или-беги}, стимулируя кору головного мозга через слова, фразы, одежду, статус,
помещение, предметы, аксессуары.\\[0.4em]
\textbf{Методы:}
\begin{itemize}
  \item Вешая на вас ярлыки
  \item Угрожая
  \item Намекая на угрозу, долг, неосознанные предположения
  \item Внушая: \kav{ты есть ярлык, и как следствие ты испытываешь вот это}
  \item Говоря строгим тоном, чтобы поднять ваш уровень кортизола~--- заставить бояться
\end{itemize}
\end{mechbox}

Например: \kav{я же вижу, что ты нервничаешь}~--- злоумышленник намеренно пытается
вызвать эту семантическую реакцию, заставить вас говорить то, о чём вы пожалеете.

Цель: \textbf{отключить префронтальную кору}:
\begin{enumerate}
  \item Дать гипофизу выплеснуть кортизол
  \item Перевести амигдалу в первое действие
  \item Не дать сигналу доходить до коры
  \item Перекоммутировать сигналы так, чтобы звук человека распознавался как угроза
\end{enumerate}

Как следствие~--- у вас не остаётся возможности провести адекватный анализ.

\begin{scheme}{Перехват коры: быстрый путь против медленного}{9cm}
\begin{tikzpicture}
  \node[cy, text width=6.2cm] (tr) at (0,0) {\textbf{Триггер}\\слово \textbullet\ тон \textbullet\ статус \textbullet\ ярлык};
  \node[cr, text width=4.6cm] (am) at (-4.2,-2.5) {\textbf{Таламус $\to$ Амигдала}\\быстрый путь, миллисекунды};
  \node[cb, text width=4.6cm] (co) at (4.2,-2.5) {\textbf{Таламус $\to$ Кора}\\медленный путь};
  \node[cr, text width=4.6cm] (ko) at (-4.2,-5.0) {\textbf{Кортизол}\\\kav{бей или беги}};
  \node[cw, dashed, text width=4.6cm] (an) at (4.2,-5.0) {\textbf{Анализ}\\не успевает~--- отключён};
  \node[cr, text width=8.4cm, line width=1.6pt] (ea) at (0,-7.4) {\termr{E.A.Alert}\\\textbf{реакция без анализа}};
  \draw[arrr] (tr.south west) ++(0.6,0) -- (am.north);
  \draw[arrd] (tr.south east) ++(-0.6,0) -- (co.north);
  \draw[arrr] (am) -- (ko);
  \draw[arrd] (co) -- (an);
  \draw[arrr] (ko.south) -- ++(0,-0.5) -| ([xshift=-2cm]ea.north);
  \draw[arrd] (an.south) -- ++(0,-0.5) -| ([xshift=2cm]ea.north);
  \node[font=\footnotesize\itshape, text=sce-red] at (0,-3.7) {кора получает сигнал позже};
  \node[pin] at (tr.north west) {}; \node[pin] at (am.north west) {}; \node[pin] at (co.north west) {};
  \node[pin] at (ko.north west) {}; \node[pin] at (an.north west) {}; \node[pin] at (ea.north west) {};
  \boardbg
\end{tikzpicture}
\end{scheme}


\begin{warnbox}{Руминация~--- амигдала берёт контроль}
Вспомнив событие/диалог/перепалку, вы можете опять перейти в режим \kav{бей-или-беги}~---
и это просто из-за \kav{образа} в голове.\\[0.4em]
Вы циклично прокручиваете диалоги: \kav{как лучше бы мне ответить, да так чтобы пожалел}.\\
У вас снова поднялся кортизол. 5 минут, 1 триггер. Вы истощены.\\[0.4em]
\textbf{Скажите себе:} \kav{Это не анализ. Стоп. Снова руминация.\\
Амигдала взяла контроль над моими семантическими реакциями на это событие.\\
Как адекватно отреагировать?}
\end{warnbox}

\subsection{Деструктивное влияние алкоголя. ИФД заразно. Влияние больных на детей и взрослых}

С такой проблемой можно в буквальном смысле \kav{заразиться} ИФД от человека, и даже
\kav{заразить} других. Алкоголь нарушает микробиом кишечника~--- надлежащее питание к
нейронам не поступает. Когда амигдала многократно переводит вас в режим \kav{бей-или-беги},
выделяя кортизол, вы не можете адекватно оценивать.

\begin{warnbox}{Никогда не оставляйте ребёнка с человеком, больным ИФД}
Иначе он переймёт его болезни. Вечные нападки доведут его до суицидальных мыслей.\\[0.3em]
Ребёнок в режиме \kav{бей-или-беги} видит мир крайне опасным~--- места кажутся
пустыми, страшными, тёмными, недоверительными по отношению к людям.\\[0.3em]
ИФД + алкоголь~--- явно ухудшает проблему. Со временем начнутся крики, семейные разбои.
\end{warnbox}

\subsection{Война оправдана гормоном, а не логикой}

Говоря о переговорах о мире~--- вы не сможете договориться, если что-то переводит
вас и вашего собеседника в режим \kav{бей-или-беги}. Амигдала диктует семантические реакции.

Чтобы побудить людей воевать~--- их сначала переводят в режим \kav{бей-или-беги},
получив контроль над амигдалой, диктуют нужные деструктивные семантические реакции:
через радио, ТВ, музыку, символику, плакаты.

\begin{warnbox}{Самое немыслимое}
Детей в детсадах и школах приучают к этим неадекватным семантическим реакциям~---
по умолчанию выставляя предустановку о нормальности убийств за королей,\\
вырабатывая агрессию на образы, знаки, символы отличия одной группы от другой.\\[0.3em]
Это~--- серьёзная опасность.
\end{warnbox}

% ============================================================
\section{Инъекция семантических реакций}
% ============================================================

\begin{mechbox}{Определение}
\texttt{Semantic-Reaction-Injection [SRI]}\\
\texttt{Инъекция Семантических Реакций [ИСР]}\\[0.4em]
Намеренное или случайное введение карт-ярлыков в оценочную систему человека,\\
изменяющее его семантические реакции в обход сознательного анализа.
\end{mechbox}

% ============================================================
\section{Модули}
% ============================================================

\begin{mechbox}{\termr{SRI.DS}~--- S.R.I. Detection System~--- Система обнаружения вторжений}
S.R.I. Detection System~--- система обнаружения вторжений, анализирующая семантические
реакции индивидуумов и их системные высказывания на предмет подозрительной словесной
активности, выявления опасных тенденций в оценках, прогнозирования инцидентов.
\end{mechbox}

\begin{mechbox}{\texttt{САРТСНА}}
Тест, используемый для различения человека и логических машин, применяемый для
предотвращения автоматических неадекватных карт и других действий через оценку.
\end{mechbox}

% ============================================================
\section{Атаки}
% ============================================================

Вирусы используются с целью подмены ваших ответных семантических реакций на неестественный,
неадекватный порядок процесса абстрагирования.

\begin{mechbox}{\termr{SM}~--- Semantic Malware~--- Семантическая малварь}
Вредоносный язык для нанесения ущерба индивидуумам. Виды включают: интенсиональные
карты, не сверенные с территорией; жестокий язык; религиозные вирусы; карты-трояны,
содержащие в себе далеко идущие неосознанные предположения.
\end{mechbox}

\begin{mechbox}{\termr{MAP Spoofing} / [MPS]~--- MAP Spoofing~--- Отравление картой (подмена карты)}
Атака, при которой злоумышленник рисует ложные карты-сообщения, давая неадекватные
термины, обобщающие выводы-ярлыки для изменения семантических реакций и наших
неосознанных оценок~--- перекоммутируя активации нейронных сетей.
\end{mechbox}

\begin{mechbox}{\termr{BruF}~--- Brute Force~--- Брутфорс}
Метод не\-/осознанной атаки, при которой злоумышленник перебирает все возможные
комбинации слов, фраз до тех пор, пока не найдёт ту, которая вас спровоцирует на
неадекватную ответную семантическую реакцию, чтобы потом сказать вам, что вы сами начали.
\end{mechbox}

\begin{mechbox}{\termr{SREx}~--- S.R. Exploit~--- Эксплойт семантических реакций}
Языковая подпрограмма, гормон-код, карта, оценки, образы, аналогии, использующие
уязвимости в процессе абстрагирования-оценки для выполнения несанкционированных действий,
с полным перехватом управления над оценкой.
\end{mechbox}

\begin{mechbox}{\termr{SR.Script}~--- S.R. Scripting~--- Скриптинг семантических реакций}
В контексте выступлений~--- процесс написания тезисов для речи с целью вызвать нужные
ответные семантические реакции. Получатель интерпретирует звук и составляет карты в голове,
которые мало соответствуют реальной территории.
\end{mechbox}

\begin{mechbox}{\termr{SRI.S}~--- S.R. Injection by Status~--- Инъекция через статус}
Определение вас по статусу, по вашему жизненному положению, и инъекция оценок,
\kav{правильных} взглядов. Сначала на вас вешают ярлык~--- потом вставляют нужную
самооценку-карту, делая из вас неуверенного в себе индивида.
\end{mechbox}

\begin{mechbox}{\termr{SRI-Priority}~--- S.R. Injection by Priority~--- Инъекция через авторитет}
Эксплуатирует чувство вины. Строгий тон, нахмуренное выражение, словесное указание
\kav{своего места}~--- автоматически переводит в режим Emygdala Alert.\\
Вы отдаёте приоритет не своим оценкам~--- просто потому что их проговорили строго.
\end{mechbox}

\begin{mechbox}{\termr{SRI-DA}~--- S.R.I. Dictionary Attack~--- Словарная атака}
Атака с использованием предварительно собранного списка слов (словаря) для вызова
неадекватных семантических реакций. Перевод индивидуума в Emygdala Alert.
\end{mechbox}

\begin{mechbox}{\termr{MTS}~--- Map Terminology Spoofing~--- Подмена терминологии карты}
Процесс, при котором злоумышленник подменяет терминологию карты для скрытия своего
личного интереса. Передаётся карта с неверной терминологией под видом объективной.
\end{mechbox}

\begin{mechbox}{\termr{GRB}~--- Gormon S.R. Backdoor~--- Гормональный бэкдор семантических реакций}
Обусловленный языковой средой с рождения гормональный программный языковой код~---
механизмы, позволяющие обходить обычные средства надлежащего порядка абстрагирования\\
и получать несанкционированный доступ к системе.
\end{mechbox}

\begin{mechbox}{\termr{CSRF}~--- Cross-S.R. Request Forgery~--- Подделка запроса между реакциями}
Атака, при которой злоумышленник заставляет индивидуума выполнить нежелательное действие
под якобы чужую ответственность. Используя семантические инъекции, человек потом
использует специальные фразы, чтобы избавиться от личной ответственности.
\end{mechbox}

\begin{mechbox}{\termr{DDoP}~--- Distributed Denial of People~--- Распределённый отказ людям (перегрузка)}
Большое количество запросов на достижение \kav{успеха} при их явной невозможности в данное
время~--- ограничение ресурсов в еде, воде и сне. Перестимуляция словами. Перевешивание ярлыками.
\end{mechbox}

\begin{mechbox}{\termr{CRS}~--- Cross ReEvulation S.R. Scripting~--- Перекрёстный скриптинг оценок}
Атака, при которой злоумышленник вставляет вредоносный скрипт (неосознанные предположения,
оценки) в оценку, которую принимают как данное другие индивидуумы~---
без каких-либо свидетельств и проверок.
\end{mechbox}

\begin{mechbox}{\termr{MITM}~--- Man-in-the-Middle~--- Человек посередине}
Атака, при которой человек умышленно и/или неосознанно перехватывает и изменяет в
процессе коммуникации оценки между двумя сторонами без их ведома.\\
Реализуется через атаку Emygdala Alert.
\end{mechbox}

% ============================================================
\section{Общие ошибки передачи данных при коммуникации}
% ============================================================

Коммуникация между людьми напоминает игру в сломанный телефон.
Каждая передача от человека к человеку искажает первичные данные.

\begin{mechbox}{\termr{EDC}~--- Extentional Data Corruption~--- Коррупция экстенсиональных данных}
Процесс преобразования экстенсиональных данных в интенсиональную неизмеримую форму,
недоступную для понимания без специального обучения наблюдению, задаванию вопросов,
выяснению фактов, экспериментов~--- обеспечивающий недоступность экстенсиональных данных
при их передаче в обучении, книгах, школах.
\end{mechbox}

\begin{mechbox}{\termr{AS.RL}~--- Abstract Semantic-Reaction Layer~--- Абстрактный слой семантических реакций}
Протоколы \kav{шифрования} в документах, договорах, в политике якобы конфиденциальности
(aka мелкий шрифт)~--- обеспечивающие общую интенсиональность в ориентировании,
ломающие передачу экстенсиональных данных. Словечники-инъекторы без физических
свидетельств и измерений.
\end{mechbox}

% ============================================================
\section{Аберрантное поведение}
% ============================================================

Поведение, возникающее как следствие накопленных семантических атак: руминация, агрессия,
тревога, неспособность к диалогу, избегание, абстрактные страхи. Не признак слабости~---
признак отсутствия инструментов для работы с оценочной средой.

% ============================================================
\section{Сверка карт и территории}
% ============================================================

\begin{mechbox}{\termr{Rootkit}~--- Rootkit (Hidden conclusions)~--- Скрытые предустановки}
Неосознанные факторы внешней/внутренней среды, скрытые и предназначенные для
обеспечения постоянного доступа. Руткиты возможно изменяют оценки в связи с
физиологическими процессами, делая их трудно обнаруживаемыми стандартными
аристотелевыми методами.
\end{mechbox}

Сверка карты и территории~--- единственное средство защиты:
\begin{itemize}
  \item Задать вопрос: \textit{«это факт или оценка?»}
  \item Снизить уровень абстракции: \textit{«что именно произошло?»}
  \item Допустить неопределённость: \textit{«я не знаю достаточно»}
  \item Не прыгать к выводам~--- \texttt{No jump to conclusions!}
\end{itemize}

% ============================================================
%  ЧАСТЬ II — Вирусология оценок: Школа
% ============================================================

\newpage
\begin{center}
\vspace*{1cm}
\textcolor{sce-gold}{\rule{0.5\textwidth}{0.4pt}}\\[0.8em]
{\montserrat\Large\bfseries\addfontfeatures{LetterSpace=3}\color{sce-black}
ЧАСТЬ II}\\[0.4em]
{\montserrat\large\addfontfeatures{LetterSpace=5}\color{sce-grey}
В И Р У С О Л О Г И Я}\\[0.2em]
{\montserrat\normalsize\addfontfeatures{LetterSpace=3}\color{sce-grey}
оценок~--- Школа}\\[0.4em]
{\small\itshape\color{sce-grey}
Дополнение к «Вирусологии семантических реакций» v\,0.1}\\[0.8em]
\textcolor{sce-gold}{\rule{0.5\textwidth}{0.4pt}}
\vspace{1cm}
\end{center}

% ============================================================
\section{Введение~--- Школа как среда передачи}
% ============================================================

Школа~--- это не нейтральная среда.\\
Это~--- \textbf{первая массовая точка передачи семантических вирусов}.

Ребёнок приходит туда с открытой системой абстрагирования.
Он ещё не имеет устоявшихся карт. Он~--- \textit{наблюдатель}.
Школа за~12 лет систематически заменяет наблюдение~--- \textbf{ярлыком}.

\begin{warnbox}{Ключевое наблюдение}
Школа~--- не обязательно злой умысел.\\
Это~--- система, в которой \textbf{учителя сами заражены}, и передают то, что имеют.\\[0.3em]
\termr{GRB} → крик учителя → \termr{GRB} ученика → крик его детей.\\[0.3em]
Цепочка. Не заговор. Вирус.
\end{warnbox}

% ============================================================
\section{Система оценок~--- \termr{MAP Spoofing} в промышленном масштабе}
% ============================================================

\subsection{Что такое оценка в школе}

Оценка в школе~--- это число от~1 до~5 (или буква, или процент).
Это~--- \textbf{карта}. Она описывает \textit{что-то} о поведении ребёнка в конкретный момент.

Но школьная система передаёт её иначе:

\begin{mechbox}{Как передаётся оценка~--- цепочка абстракций}
\textbf{Уровень 0 (событие):} ребёнок не смог пересказать текст в 14:30 во вторник\\
\textbf{Уровень 1 (оценка):} \texttt{«2»} в журнале\\
\textbf{Уровень 2 (обобщение):} \texttt{«плохо учится»}\\
\textbf{Уровень 3 (личность):} \texttt{«лентяй», «тупой», «неспособный»}\\
\textbf{Уровень 4 (прогноз):} \texttt{«ничего не добьётся»}\\[0.4em]
Каждый уровень выше~--- \textbf{дальше от территории}.\\
Школа работает на уровне~3--4. Это~--- \termr{MAP Spoofing}.
\end{mechbox}

\begin{scheme}{Лестница абстракций: от события к ярлыку}{9cm}
\begin{tikzpicture}
  \node[cg, text width=2.65cm, anchor=north west] (a0) at (0,0)      {\ttl УРОВЕНЬ 0 \\[1pt] \textbf{Событие}\\не смог пересказать текст};
  \node[cb, text width=2.65cm, anchor=north west] (a1) at (3.1,-0.9) {\ttl УРОВЕНЬ 1 \\[1pt] \textbf{Оценка}\\\kav{2} в журнале};
  \node[cy, text width=2.65cm, anchor=north west] (a2) at (6.2,-1.8) {\ttl УРОВЕНЬ 2 \\[1pt] \textbf{Обобщение}\\\kav{плохо учится}};
  \node[cr, text width=2.65cm, anchor=north west] (a3) at (9.3,-2.7) {\ttl УРОВЕНЬ 3 \\[1pt] \textbf{Личность}\\\kav{лентяй}, \kav{тупой}};
  \node[cr, text width=2.65cm, anchor=north west, line width=1.6pt] (a4) at (12.4,-3.6) {\ttl УРОВЕНЬ 4 \\[1pt] \textbf{Прогноз}\\\kav{ничего не добьётся}};
  \foreach \a/\b in {a0/a1,a1/a2,a2/a3,a3/a4}
    \draw[thread] (\a.east) ++(0,-0.2) to[out=0,in=180] (\b.west);
  \foreach \n in {a0,a1,a2,a3,a4} \node[pin] at (\n.north west) {};
  \draw[arr] (0,-7.0) -- node[below=2pt,font=\footnotesize\itshape] {расстояние от территории растёт} (15.4,-7.0);
  \node[font=\footnotesize\bfseries, text=sce-green, anchor=south west] at (0,-6.95) {ТЕРРИТОРИЯ};
  \node[font=\footnotesize\bfseries, text=sce-red, anchor=south east] at (15.4,-6.95) {КАРТА};
  \draw[decorate,decoration={brace,amplitude=5pt,mirror}, sce-red, line width=1.2pt]
     (9.3,-6.4) -- (15.4,-6.4) node[midway,below=6pt, font=\footnotesize\bfseries, text=sce-red] {};
  \node[font=\footnotesize\bfseries, text=sce-red] at (12.35,-8.1) {школа работает здесь \termr{MAP Spoofing}};
  \boardbg
\end{tikzpicture}
\end{scheme}


\subsection{Почему это \termr{MAP Spoofing}}

\termr{MAP Spoofing}~--- атака, при которой система даёт \textbf{ложную карту-ярлык}
вместо описания реального события. Школьная оценка как ярлык:
\begin{itemize}
  \item Не описывает \textit{что именно} не получилось
  \item Не объясняет \textit{почему}
  \item Не содержит \textit{инструмента} для исправления
  \item Навешивается на \textbf{личность} ребёнка
\end{itemize}

\begin{dialogbox}{Пример~--- реальная цепочка}
Учитель (вслух, перед классом):\\
\texttt{«Опять двойка. Совсем не стараешься.»}\\[0.3em]
Что произошло: ребёнок не смог пересказать текст.\\
Что было сказано: \texttt{[ты] = [не стараешься]}~--- оценка личности.\\
Что осталось у ребёнка: карта себя как \texttt{«того, кто не старается»}.\\[0.3em]
Территория: нет метода работы с текстом.\\
Карта: \textit{лентяй}.
\end{dialogbox}

% ============================================================
\section{\termr{SRI} в классе~--- Инъекции оценок}
% ============================================================

\subsection{Публичное унижение как \termr{E.A.Alert}}

Когда учитель зачитывает оценки вслух перед классом~---
это \textbf{\termr{E.A.Alert}} для ребёнка с низкой оценкой.

\begin{mechbox}{Механизм срабатывания}
1. Учитель называет имя + двойку при всех.\\
2. Класс смотрит. Смех. Или тишина~--- что хуже.\\
3. Амигдала: угроза социального статуса = угроза выживания.\\
4. Кортизол поднимается.\\
5. Кора отключается. Анализ невозможен.\\
6. Ребёнок в режиме \texttt{«бей или беги»}.\\[0.3em]
В этом состоянии он \textbf{не может учиться}. Но урок продолжается.
\end{mechbox}

\subsection{Повторяющийся \termr{GRB}~--- формирование рефлекса}

Если \termr{E.A.Alert} в классе повторяется регулярно~---
у ребёнка формируется \termr{GRB}:

\begin{mechbox}{\termr{GRB} в школьном контексте}
\textbf{Триггер:} ситуация проверки знаний (контрольная, вопрос у доски, пересказ)\\
\textbf{Гормональный ответ:} кортизол → амигдала → режим угрозы\\
\textbf{Поведение:} ступор, забывание, заикание, агрессия, уход в себя\\
\textbf{Карта ребёнка о себе:} \texttt{«я не могу»}\\[0.3em]
Позже: взрослый человек в любой ситуации проверки~---\\
получает тот же кортизол. \textbf{Хотя угрозы нет.}\\[0.3em]
Это~--- \termr{GRB}. Школьный бэкдор. Пожизненный.
\end{mechbox}

\subsection{\termr{SRI-Priority}~--- строгий тон как инструмент}

Учитель с авторитетным тоном~--- это \termr{SRI-Priority}:
ребёнок принимает карту учителя \textbf{без проверки}~--- даже если карта ложная.
Результат~--- \textbf{интенсиональная ориентация}: ребёнок живёт не по факту,
а по ярлыку.

% ============================================================
\section{Вирусология оценок~--- классификация школьных атак}
% ============================================================

\begin{mechbox}{Таблица~--- Атаки в школьной среде}
\begin{tabularx}{\linewidth}{@{} l X @{}}
\toprule
\textbf{Код} & \textbf{Описание атаки в школе} \\
\midrule
\termr{MAP Spoofing} &
  Оценка как ярлык личности вместо описания конкретного действия. \\[0.3em]
\termr{E.A.Alert} &
  Публичное объявление оценки. Класс как аудитория угрозы. \\[0.3em]
\termr{GRB} &
  Выработка гормонального бэкдора: ситуация проверки = угроза. Пожизненный рефлекс. \\[0.3em]
\termr{SRI-Priority} &
  Строгий тон учителя как авторитет. Интенсиональная ориентация. \\[0.3em]
\termr{SRI} &
  Инъекция: \textit{«ты должен», «все нормальные дети умеют», «тебе стыдно должно быть»}. \\[0.3em]
\termr{DDoP} &
  Перегрузка: домашнее задание + контрольные + требования + оценки + родители. \\[0.3em]
\termr{EDC} &
  Зубрёжка: наблюдение и понимание заменены формулировками из учебника. \\[0.3em]
\termr{Rootkit} &
  Скрытая предустановка: \textit{«хорошая оценка = ты хороший человек»}. \\[0.3em]
\termr{CRS} &
  \textit{«Петя умеет, а ты нет»}. Чужая оценка как эталон. \\
\bottomrule
\end{tabularx}
\end{mechbox}

% ============================================================
\section{Лень не существует~--- \termr{MAP Spoofing} имя которому «лень»}
% ============================================================

\begin{warnbox}{Центральный тезис}
Ярлык \texttt{[ ЛЕНЬ ]}~--- это \textbf{\termr{MAP Spoofing}}.\\
Он описывает не территорию, а \textit{отсутствие карты у наблюдателя}.\\[0.3em]
Когда учитель (или родитель) говорит «ленится»~---\\
это означает: \textit{я не вижу реальной причины. Я использую ярлык.}
\end{warnbox}

За ярлыком \texttt{[ ЛЕНЬ ]} в школьном контексте~--- всегда одно из:

\begin{factbox}{Что скрыто за ярлыком~--- территория}
\begin{itemize}
  \item \textbf{Нет метода}~--- ребёнок не умеет работать с этим типом задачи. Это \termr{EDC}.
  \item \textbf{Нет смысла}~--- связь между задачей и жизнью не объяснена. Мозг экономит ресурс.
  \item \textbf{\termr{GRB} активен}~--- кортизол. В режиме угрозы обучение биологически заблокировано.
  \item \textbf{Истощение~--- \termr{DDoP}}~--- 6 уроков + домашнее + кружок + тревога дома.
  \item \textbf{Скрытый протест}~--- отказ от навязанного. Сигнал о несоответствии.
  \item \textbf{\termr{FEAR OF ERROR}}~--- страх ошибиться парализует действие.
\end{itemize}
\end{factbox}

\begin{scheme}{Что стоит за ярлыком [ ЛЕНЬ ]: причина и инструмент}{10cm}
\begin{tikzpicture}
  \node[cr, text width=2.9cm, line width=1.6pt] (lz) at (0.6,-3.3) {\textbf{[ ЛЕНЬ ]}\\ярлык\\\termr{MAP Spoofing}};
  \foreach \i/\y/\c/\cause/\tool in {
     1/0/cb/{Нет метода}/{показать шаг: \kav{вот как делаю я}},
     2/-1.35/cb/{Нет смысла}/{показать связь с жизнью},
     3/-2.7/cr/{\termr{GRB} активен}/{пауза, снять угрозу},
     4/-4.05/cy/{Истощение \termr{DDoP}}/{убрать лишнее, дать ресурс},
     5/-5.4/cy/{Скрытый протест}/{спросить, что не так},
     6/-6.75/cr/{\termr{FEAR OF ERROR}}/{переопределить ошибку}}
  {
    \node[\c, text width=4.0cm] (c\i) at (6.3,\y) {\textbf{\cause}};
    \node[cg, text width=4.9cm] (t\i) at (12.2,\y) {\tool};
    \draw[thread] (lz.east) to[out=0,in=180] (c\i.west);
    \draw[arrg] (c\i.east) -- (t\i.west);
    \node[pin] at (c\i.north west) {};
  }
  \node[font=\ttl, text=sce-grey] at (6.3,1.0) {ТЕРРИТОРИЯ: реальная причина};
  \node[font=\ttl, text=sce-green] at (12.2,1.0) {ИНСТРУМЕНТ ВМЕСТО ЯРЛЫКА};
  \boardbg
\end{tikzpicture}
\end{scheme}


\subsection{Цепочка передачи ярлыка}

\begin{center}
\begin{tabular}{ccccc}
Школа & $\to$ & \termr{GRB} учителя & $\to$ & Ярлык «лентяй» \\
$\downarrow$ & & & & $\downarrow$ \\
Нет метода & & & & \termr{GRB} у ребёнка \\
$\downarrow$ & & & & $\downarrow$ \\
Ребёнок идёт домой & $\to$ & Мать видит двойку & $\to$ & \termr{GRB} матери \\
& & $\downarrow$ & & \\
& & Крик & $\to$ & \termr{FEAR OF ERROR} \\
& & $\downarrow$ & & \\
& & \multicolumn{3}{l}{\texttt{[ ЛЕНЬ ]}~--- ярлык закрепляется в карте о себе} \\
\end{tabular}
\end{center}

Ни учитель, ни мать не осознают цепочку. Они \textbf{передают то, что сами получили}.\\
Это~--- \textbf{вирус}.

\begin{scheme}{Цепочка заражения: как вирус передаётся по кругу}{10cm}
\begin{tikzpicture}
  \node[cy, text width=3.6cm] (n1) at (0,3.3)    {\textbf{Школа}\\оценка как ярлык};
  \node[cr, text width=3.6cm] (n2) at (5.0,1.7)  {\textbf{Учитель}\\\termr{GRB}: давление, тон};
  \node[cr, text width=3.6cm] (n3) at (5.0,-1.7) {\textbf{Ребёнок}\\\termr{E.A.Alert} в классе};
  \node[cr, text width=3.6cm] (n4) at (0,-3.3)   {\textbf{Родитель дома}\\\termr{GRB}: двойка в дневнике};
  \node[cr, text width=3.6cm] (n5) at (-5.0,-1.7){\textbf{Крик}\\вместо вопроса};
  \node[cr, text width=3.6cm, line width=1.6pt] (n6) at (-5.0,1.7) {\textbf{Ребёнок}\\\termr{FEAR OF ERROR}, ярлык \kav{лень}};
  \draw[arrr] (n1) -- (n2);
  \draw[arrr] (n2) -- (n3);
  \draw[arrr] (n3) -- (n4);
  \draw[arrr] (n4) -- (n5);
  \draw[arrr] (n5) -- (n6);
  \draw[arrr] (n6) -- node[above left=-1pt,font=\footnotesize\itshape,text=sce-red,pos=0.45]{следующее поколение} (n1);
  \node[font=\montserrat\bfseries\normalsize, text=sce-red, align=center] at (0,0) {ВИРУС\\НЕ УМИРАЕТ};
  \foreach \n in {n1,n2,n3,n4,n5,n6} \node[pin] at (\n.north west) {};
  \boardbg
\end{tikzpicture}
\end{scheme}


% ============================================================
\section{Зубрёжка~--- \termr{EDC} как системная атака}
% ============================================================

\subsection{Механизм EDC в учёбе}

\begin{mechbox}{Нормальный процесс vs. школьная зубрёжка}
\textbf{Нормальный процесс абстрагирования:}\\
Событие (наблюдение) → Описание → Обобщение → Принцип\\[0.4em]
\textbf{Школьная зубрёжка:}\\
Формулировка из учебника → Запомнить → Воспроизвести → Оценка\\[0.4em]
Наблюдение~--- пропущено. Понимание~--- необязательно.\\
Связь с реальностью~--- отсутствует.
\end{mechbox}

\subsection{Пример: стихотворение}

\begin{dialogbox}{Реальная ситуация (флэшбэк)}
Задание: выучить стихотворение наизусть.\\
Цель по программе: развитие памяти и \kav{приобщение к культуре}.\\[0.4em]
Что происходит в реальности:\\
Ребёнок не понимает текст. Слова~--- абстрактные.\\
Нет метода: как работать с непонятным текстом?\\
Результат: зубрёжка без смысла. Или~--- отказ. Или~--- \termr{GRB}.\\[0.3em]
Ярлык: \texttt{[ ЛЕНЬ ]}.\\
Территория: \textit{нет инструмента работы с непонятным}.
\end{dialogbox}

\subsection{Инструмент вместо зубрёжки}

\begin{factbox}{Четыре шага вместо \kav{учи наизусть}}
1. \textbf{Прочитай первые две строки. Скажи своими словами~--- о чём?}\\
   (Экстенсиональное наблюдение: что я понял?)\\[0.3em]
2. \textbf{Найди одно слово, которое непонятно.}\\
   (Конкретизация: что именно мешает?)\\[0.3em]
3. \textbf{Придумай образ к каждой строке.}\\
   (Связь с опытом: карта ↔ территория)\\[0.3em]
4. \textbf{Расскажи~--- не повторяй, а расскажи.}\\
   (Проверка понимания, а не памяти)\\[0.3em]
Это занимает те же 20 минут.\\
Но ребёнок получает \textbf{метод}, а не рефлекс страха.
\end{factbox}

% ============================================================
\section{Патриотизм в школе~--- \termr{SR.Script} и \termr{Rootkit}}
% ============================================================

\subsection{Механизм формирования}

\begin{mechbox}{Компоненты \termr{SR.Script} в школе}
\textbf{Символика:} флаг, герб, гимн~--- якоря для \termr{E.A.Alert}\\
\textbf{Нарратив:} история как набор правильных карт\\
\textbf{Ритуал:} линейка, минута молчания, \kav{Россия~--- великая страна}\\
\textbf{Оценка:} правильный ответ о \kav{родине} = хорошая оценка\\
\textbf{Давление:} неправильный ответ = моральное осуждение коллектива\\[0.3em]
Результат: \termr{Rootkit}~--- скрытая предустановка:\\
\textit{«моя страна всегда права»}~--- работает без осознания.
\end{mechbox}

\subsection{Почему это \termr{Rootkit}}

\termr{Rootkit}~--- скрытые предположения, которые \textbf{не осознаются} и
предоставляют \textbf{постоянный доступ} к оценкам:
\begin{itemize}
  \item Не подвергается проверке~--- это \textit{само собой разумеющееся}
  \item Активируется автоматически при триггере (флаг, гимн, \kav{враг})
  \item Блокирует критический анализ: \textit{«это~--- предательство»}
  \item Передаётся следующему поколению через те же ритуалы
\end{itemize}

\subsection{Агрессия родителей~--- производная школьной среды}

Когда школа требует патриотического поведения + оценки + результат~---
она передаёт \textbf{давление} родителям:

\begin{enumerate}
  \item Школа: \textit{«ваш ребёнок не выполняет требования»}
  \item Родитель: \termr{GRB} активируется~--- стыд, страх, злость
  \item Дома: крик, давление на ребёнка~--- \textit{«почему двойка?»}
  \item Ребёнок: \termr{E.A.Alert} + \termr{GRB} → \termr{FEAR OF ERROR}
  \item Результат: ребёнок не учится лучше. Он учится \textbf{бояться}.
\end{enumerate}

\begin{dialogbox}{Ключевое наблюдение}
Родитель не хочет причинять вред.\\
Родитель~--- сам заражён \termr{GRB} от \textit{своей} школы.\\
Крик~--- это не жестокость. Это~--- \textbf{единственный инструмент},\\
который был передан ему.\\[0.3em]
Школа не дала методов~--- она дала \textbf{давление + оценку}.\\
Родитель транслирует то, что получил.\\[0.3em]
Вирус~--- передался.
\end{dialogbox}

% ============================================================
\section{Последствия~--- карта себя}
% ============================================================

Ребёнок, прошедший 10--12 лет в такой системе, выходит с набором карт:

\begin{warnbox}{Карты, которые выдаёт школьная система}
\texttt{[я = моя оценка]}~--- ценность личности определяется числом\\
\texttt{[ошибка = опасность]}~--- \termr{FEAR OF ERROR} как базовая установка\\
\texttt{[не знаю = стыд]}~--- незнание~--- это не начало, это~--- провал\\
\texttt{[авторитет прав]}~--- карта принята без проверки~--- \termr{SRI-Priority}\\
\texttt{[я лентяй/тупой/неспособный]}~--- ярлык принят как факт\\
\texttt{[думать опасно]}~--- правильный ответ важнее честного наблюдения
\end{warnbox}

\subsection{Что не передала школа}

\begin{factbox}{Что должно было быть передано}
\begin{itemize}
  \item Как работать с непонятным текстом (метод, а не зубрёжка)
  \item Как задавать вопросы (экстенсиональная ориентация)
  \item Как замечать, что ты не знаешь~--- и считать это началом, не провалом
  \item Как проверять карту~--- \textit{«это факт или оценка?»}
  \item Как справляться с ошибкой~--- \textit{«что именно не получилось?»}
  \item Как распознавать \termr{E.A.Alert} в себе и делать паузу
\end{itemize}
\textbf{Это~--- Общая Семантика.}\\
Её не преподают. Потому что система сама работает на \termr{MAP Spoofing}.
\end{factbox}

% ============================================================
\section{Антивирус~--- что работает}
% ============================================================

\subsection{Признание уровня абстракции}

\begin{mechbox}{Вопросы для проверки карты}
\textbf{Это факт или оценка?}\\
\textit{«Он лентяй»}~--- это оценка. \textit{«Он не сдал три задания»}~--- ближе к факту.\\[0.3em]
\textbf{Что именно произошло?}\\
Не \textit{«плохо учится»}, а: \textit{«в задаче на дроби пропускает шаг сокращения»}.\\[0.3em]
\textbf{Какой инструмент нужен?}\\
Не \textit{«старайся больше»}, а: \textit{«давай разберём этот шаг по-другому»}.
\end{mechbox}

\subsection{Пауза перед \termr{GRB}}

\termr{GRB}~--- гормональный. Его нельзя отключить мыслью. Но можно~--- \textbf{создать паузу}.

\begin{factbox}{Метод паузы}
1. Заметить: \textit{«у меня сейчас кортизол. Что его вызвало?»}\\
2. Назвать триггер: \textit{«двойка в дневнике»}~--- это событие, не катастрофа.\\
3. Отделить событие от оценки:\\
   \textit{«Двойка»}~$\neq$~\textit{«плохой ребёнок»}~$\neq$~\textit{«я плохой родитель»}\\
4. Задать вопрос вместо крика:\\
   \textit{«Что именно не получилось?»}
\end{factbox}

\begin{scheme}{Антивирус: пауза между триггером и реакцией}{5cm}
\begin{tikzpicture}
  \node[cr, text width=2.9cm] (a) at (0,0)    {\textbf{Триггер}\\\termr{GRB} запущен: кортизол};
  \node[cy, text width=2.9cm] (b) at (3.4,0)  {\textbf{1. Заметить}\\\kav{у меня сейчас кортизол}};
  \node[cb, text width=2.9cm] (c) at (6.8,0)  {\textbf{2. Назвать}\\что именно запустило};
  \node[cb, text width=2.9cm] (d) at (10.2,0) {\textbf{3. Отделить}\\событие от оценки};
  \node[cg, text width=2.9cm, line width=1.6pt] (e) at (13.6,0) {\textbf{4. Спросить}\\\kav{что именно не получилось?}};
  \foreach \a/\b in {a/b,b/c,c/d,d/e} \draw[arrg] (\a) -- (\b);
  \node[font=\footnotesize\bfseries, text=sce-green] at (8.5,-1.9) {пауза: здесь заканчивается рефлекс и начинается выбор};
  \draw[decorate,decoration={brace,amplitude=5pt,mirror}, sce-green, line width=1.2pt] (3.4,-1.2) -- (13.6,-1.2);
  \foreach \n in {a,b,c,d,e} \node[pin] at (\n.north west) {};
  \boardbg
\end{tikzpicture}
\end{scheme}


\subsection{Инструмент вместо ярлыка~--- три шага}

\begin{factbox}{Три шага вместо ярлыка «лень»}
\textbf{1. Нет метода → Показать шаг}\\
Не \textit{«старайся»}, а: \textit{«вот как я делаю это. Попробуй так же»}.\\[0.3em]
\textbf{2. Нет смысла → Показать связь}\\
Не \textit{«надо знать»}, а: \textit{«это нужно тогда, когда...»}.\\[0.3em]
\textbf{3. Страх ошибки → Переопределить ошибку}\\
Не \textit{«снова неправильно»}, а: \textit{«отлично. Теперь мы знаем, что именно не так»}.
\end{factbox}

% ============================================================
\section{Итог~--- Школа в терминах вирусологии}
% ============================================================

\begin{warnbox}{Итоговая карта}
\textbf{Школа как среда:}\\
Массовая точка передачи \termr{GRB}, \termr{MAP Spoofing}, \termr{EDC}, \termr{Rootkit}.\\[0.3em]
\textbf{Механизм:}\\
Оценка как ярлык личности~--- вместо описания действия.\\
Публичность как \termr{E.A.Alert}~--- вместо безопасной среды.\\
Зубрёжка как \termr{EDC}~--- вместо метода работы с непонятным.\\
Патриотизм как \termr{Rootkit}~--- вместо навыка проверки карты.\\[0.3em]
\textbf{Результат:}\\
\termr{FEAR OF ERROR}~--- пожизненный рефлекс.\\
Карта себя~--- как набор ярлыков, принятых без проверки.\\[0.3em]
\textbf{Передача:}\\
Дети → родители (давление дома).\\
Дети → следующее поколение (тот же \termr{GRB}).\\
Вирус~--- не умирает. Он передаётся.\\[0.3em]
\textbf{Антивирус:}\\
Инструмент вместо ярлыка. Пауза вместо \termr{GRB}.\\
Вопрос вместо крика. Это~--- Общая Семантика.
\end{warnbox}

% ============================================================
% ГЛОССАРИЙ
% ============================================================
\newpage
\section*{Глоссарий терминов}
\addcontentsline{toc}{section}{Глоссарий терминов}

Все коды в тексте~--- кликабельные ссылки на этот раздел.\\
Формат: \texttt{[КОД]} \textbf{English term}, затем \textit{Перевод} и описание.

\vspace{0.6em}

\termdef{E.A.Alert}
\noindent\textbf{\texttt{[E.A.Alert]}} \quad \textbf{Emygdala Alert}\\
\textit{Перевод:} Тревога амигдалы\\
Перехват управления амигдалой в обход коры. Режим \kav{бей-или-беги} запускается словом, тоном, статусом~--- без реальной физической угрозы. Анализ блокируется.

\vspace{0.5em}
\termdef{GRB}
\noindent\textbf{\texttt{[GRB]}} \quad \textbf{Gormon S.R. Backdoor}\\
\textit{Перевод:} Гормональный бэкдор семантических реакций\\
Условный гормональный рефлекс, сформированный языковой средой с детства. Триггер (слово, ситуация, символ) вызывает кортизол автоматически~--- без осознания. Пожизненный, если не обнаружен.

\vspace{0.5em}
\termdef{SRI}
\noindent\textbf{\texttt{[SRI]} / \texttt{[ИСР]}} \quad \textbf{Semantic-Reaction Injection}\\
\textit{Перевод:} Инъекция семантических реакций (ИСР)\\
Намеренное или случайное введение карт-ярлыков в оценочную систему человека, изменяющее его реакции в обход сознательного анализа.

\vspace{0.5em}
\termdef{SRI.DS}
\noindent\textbf{\texttt{[SRI.DS]} \quad \textbf{S.R.I. Detection System}}\\
\textit{Перевод:} Система обнаружения вторжений. Анализирует семантические реакции и высказывания на предмет подозрительной активности, опасных тенденций в оценках, прогноза инцидентов.

\vspace{0.5em}
\termdef{SRI-Priority}
\noindent\textbf{\texttt{[SRI-Priority]}} \quad \textbf{S.R. Injection by Priority}\\
\textit{Перевод:} Инъекция через авторитет\\
Строгий тон, статус, нахмуренный взгляд заставляют принять карту без проверки. Человек отдаёт приоритет чужой оценке~--- просто потому что она была произнесена с авторитетом.

\vspace{0.5em}
\termdef{MAP Spoofing}
\noindent\textbf{\texttt{[MAP Spoofing]} / \texttt{[MPS]}} \quad \textbf{MAP Spoofing}\\
\textit{Перевод:} Отравление картой (подмена карты)\\
Подача ложной карты-ярлыка вместо описания реального события. Классика: оценка личности вместо описания конкретного действия. Работает через \termr{EDC}~--- замену наблюдения ярлыком.

\vspace{0.5em}
\termdef{EDC}
\noindent\textbf{\texttt{[EDC]}} \quad \textbf{Extentional Data Corruption}\\
\textit{Перевод:} Коррупция экстенсиональных данных\\
Замена наблюдения (что именно произошло?) на ярлык или формулировку. Зубрёжка~--- типичный EDC: слова без понимания. Оценка личности вместо описания действия~--- тоже EDC.

\vspace{0.5em}
\termdef{Rootkit}
\noindent\textbf{\texttt{[Rootkit]}} \quad \textbf{Rootkit (Hidden conclusions)}\\
\textit{Перевод:} Скрытые предустановки\\
Неосознанные предположения, работающие в фоне и предоставляющие постоянный доступ к оценкам. Пример: \kav{хорошая оценка = хороший человек}. Не осознаётся~--- поэтому не проверяется.

\vspace{0.5em}
\termdef{DDoP}
\noindent\textbf{\texttt{[DDoP]}} \quad \textbf{Distributed Denial of People}\\
\textit{Перевод:} Распределённый отказ людям (перегрузка)\\
Перегрузка запросами~--- требований, задач, ожиданий~--- при невозможности их выполнить в данный момент. Истощает ресурс, блокирует обучение. Школьный аналог: 6 уроков + домашнее + кружок + тревога.

\vspace{0.5em}
\termdef{FEAR OF ERROR}
\noindent\textbf{\texttt{[FEAR\_OF\_ERROR: HIGH]}} \quad \textbf{FEAR\_OF\_ERROR: HIGH}\\
\textit{Перевод:} Страх ошибки: высокий\\
Рефлекс, при котором возможная ошибка блокирует действие. Формируется через повторяющийся \termr{GRB} при ситуациях проверки. Результат: лучше не делать, чем сделать неправильно.

\vspace{0.5em}
\termdef{SR.Script}
\noindent\textbf{\texttt{[SR.Script]}} \quad \textbf{S.R. Scripting}\\
\textit{Перевод:} Скриптинг семантических реакций\\
Заготовленные тезисы, символы, ритуалы, вызывающие нужные реакции у аудитории. Используется в пропаганде, рекламе, школьных линейках. Работает через \termr{E.A.Alert}.

\vspace{0.5em}
\termdef{SM}
\noindent\textbf{\texttt{[SM]}} \quad \textbf{Semantic Malware}\\
\textit{Перевод:} Семантическая малварь\\
Вредоносный язык: интенсиональные карты без проверки, жестокий язык, религиозные вирусы, карты-трояны со скрытыми предположениями.

\vspace{0.5em}
\termdef{BruF}
\noindent\textbf{\texttt{[BruF]}} \quad \textbf{Brute Force}\\
\textit{Перевод:} Брутфорс\\
Перебор провоцирующих слов/фраз до нахождения триггера~--- после чего жертва сама \kav{начала}. Неосознанная версия~--- бытовые перепалки.

\vspace{0.5em}
\termdef{SREx}
\noindent\textbf{\texttt{[SREx]}} \quad \textbf{S.R. Exploit}\\
\textit{Перевод:} Эксплойт семантических реакций\\
Языковая конструкция, использующая уязвимость в процессе абстрагирования. Полный перехват оценки без осознания атаки.

\vspace{0.5em}
\termdef{SRI.S}
\noindent\textbf{\texttt{[SRI.S]}} \quad \textbf{S.R. Injection by Status}\\
\textit{Перевод:} Инъекция через статус\\
Определение человека через его положение~--- затем вставка \kav{правильной} самооценки. Инструмент подавления через классификацию.

\vspace{0.5em}
\termdef{SRI-DA}
\noindent\textbf{\texttt{[SRI-DA]}} \quad \textbf{S.R.I. Dictionary Attack}\\
\textit{Перевод:} Словарная атака\\
Заготовленный список слов-триггеров для перевода в \termr{E.A.Alert}. Словари ненависти, оскорблений~--- типичный пример.

\vspace{0.5em}
\termdef{MTS}
\noindent\textbf{\texttt{[MTS]}} \quad \textbf{Map Terminology Spoofing}\\
\textit{Перевод:} Подмена терминологии карты\\
Подмена слов в карте для скрытия личного интереса. \kav{Война} → \kav{специальная операция}. \kav{Убийство} → \kav{нейтрализация}.

\vspace{0.5em}
\termdef{CSRF}
\noindent\textbf{\texttt{[CSRF]}} \quad \textbf{Cross-S.R. Request Forgery}\\
\textit{Перевод:} Подделка запроса между реакциями\\
Злоумышленник заставляет человека действовать~--- снимая с себя ответственность специальными фразами. \kav{Ты сам виноват}. \kav{Я тебя предупреждал}.

\vspace{0.5em}
\termdef{CRS}
\noindent\textbf{\texttt{[CRS]}} \quad \textbf{Cross ReEvulation S.R. Scripting}\\
\textit{Перевод:} Перекрёстный скриптинг оценок\\
Вредоносная оценка одного человека принимается другими без проверки. \kav{Петя умеет~--- а ты нет}. Чужая карта становится своей.

\vspace{0.5em}
\termdef{MITM}
\noindent\textbf{\texttt{[MITM]}} \quad \textbf{Man-in-the-Middle}\\
\textit{Перевод:} Человек посередине\\
Перехват и искажение оценок между двумя сторонами. Реализуется через \termr{E.A.Alert}: пока обе стороны в режиме угрозы~--- посредник управляет передачей.

\vspace{0.5em}
\termdef{AS.RL}
\noindent\textbf{\texttt{[AS.RL]}} \quad \textbf{Abstract Semantic-Reaction Layer}\\
\textit{Перевод:} Абстрактный слой семантических реакций\\
Протоколы \kav{шифрования} в документах, договорах, мелком шрифте. Обеспечивают интенсиональность~--- ломают передачу экстенсиональных данных.

% ============================================================
\vspace{3em}
\begin{center}
\textcolor{sce-gold}{\rule{0.5\textwidth}{0.4pt}}\\[0.6em]
{\montserrat\bfseries\small\addfontfeatures{LetterSpace=6}\color{sce-black}
S O C I O C Y B E R N E E R I N G \ I N C .}\\[0.2em]
{\footnotesize\color{sce-grey}
\href{https://sociocyber.site}{sociocyber.site} \quad|\quad
GPL-3.0 \quad|\quad 2026}
\end{center}

% ============================================================
%  ПРИЛОЖЕНИЕ — Решение-проблем. Рабочая тетрадь (PDF целиком)
% ============================================================
\newpage
\phantomsection
\addcontentsline{toc}{section}{Приложение: Решение-проблем. Рабочая тетрадь}
\thispagestyle{empty}
\begin{center}
\vspace*{5cm}
\textcolor{sce-gold}{\rule{0.5\textwidth}{0.6pt}}\\[1.2em]
{\montserrat\large\bfseries\addfontfeatures{LetterSpace=6}\color{sce-black}ПРИЛОЖЕНИЕ}\\[1.2em]
{\montserrat\fontsize{24}{30}\selectfont\bfseries\color{sce-black}Решение-проблем\\Рабочая тетрадь}\\[1.2em]
{\montserrat\normalsize\bfseries\color{sce-grey}Шесть шагов от проблемы к действию}\\[1.4em]
\textcolor{sce-gold}{\rule{0.5\textwidth}{0.6pt}}\\[1.4em]
{\small\color{sce-black}Инструмент вместо ярлыка: продолжение частей I и II\\
SocioCyberneering Inc.\ \textbullet\ 19 ноября 2026 г.}
\end{center}
\includepdf[pages=-,pagecommand={\thispagestyle{empty}}]{appendix/problem_solving_ru.pdf}

\end{document}
